\documentclass[11pt,a4paper]{article}
\usepackage[utf8]{inputenc}
\usepackage[T1]{fontenc}
\usepackage[french]{babel}
\usepackage{amsmath,amssymb}
\usepackage{geometry}
\usepackage{enumitem}
\usepackage{booktabs}
\usepackage{array}
\usepackage{tabularx}
\usepackage[table]{xcolor}   % <-- option "table" pour \rowcolor
\usepackage{tcolorbox}
\usepackage{tikz}
\usepackage{fontawesome5}

\tcbuselibrary{skins,breakable}

\geometry{margin=1.8cm}

% Couleurs
\definecolor{primary}{RGB}{31,78,121}
\definecolor{soft}{RGB}{230,239,247}
\definecolor{amber}{RGB}{178,91,0}          % <-- définie
\definecolor{amberbg}{RGB}{255,244,224}
\definecolor{amberborder}{RGB}{240,192,112}
\definecolor{greenbg}{RGB}{233,246,234}
\definecolor{greenborder}{RGB}{182,224,184}
\definecolor{redbg}{RGB}{253,236,236}
\definecolor{redborder}{RGB}{242,179,171}
\definecolor{green}{RGB}{46,125,50}
\definecolor{red}{RGB}{211,47,47}

% Boîtes
\newtcolorbox{idee}[1][]{
  colback=soft, colframe=primary, boxrule=1pt, arc=6pt,
  fonttitle=\bfseries\large, coltitle=white, colbacktitle=primary,
  left=8pt, right=8pt, top=6pt, bottom=6pt,
  breakable, #1
}

\newtcolorbox{retenir}[1][]{
  colback=greenbg, colframe=greenborder, boxrule=1pt, arc=6pt,
  fonttitle=\bfseries\large, coltitle=black, colbacktitle=greenborder,
  left=8pt, right=8pt, top=6pt, bottom=6pt,
  breakable, #1
}

\newtcolorbox{attention}[1][]{
  colback=redbg, colframe=redborder, boxrule=1pt, arc=6pt,
  fonttitle=\bfseries\large, coltitle=black, colbacktitle=redborder,
  left=8pt, right=8pt, top=6pt, bottom=6pt,
  breakable, #1
}

\newtcolorbox{labo}[1][]{
  colback=amberbg, colframe=amberborder, boxrule=1pt, arc=6pt,
  fonttitle=\bfseries\large, coltitle=black, colbacktitle=amberborder,
  left=8pt, right=8pt, top=6pt, bottom=6pt,
  breakable, #1
}

\title{\textcolor{primary}{\textbf{La norme ISO 5725 — L'essentiel en 5 minutes}}\\
\large Fiche de synthèse pour étudiants BTS Bioanalyses en Laboratoire de Contrôle}
\author{}
\date{}

\begin{document}

\maketitle

\vspace{-0.5cm}

\begin{idee}[title={\faLightbulb\quad En une phrase}]
L'ISO 5725, c'est le \textbf{mode d'emploi international} pour vérifier qu'une méthode de mesure au laboratoire donne des résultats \textbf{justes} et \textbf{fidèles}.
\end{idee}

\section*{\faBullseye\quad Pourquoi ça existe ?}

\begin{itemize}[leftmargin=1.2cm, itemsep=2pt]
    \item[\textcolor{primary}{\faCheck}] Pour savoir si une méthode de mesure est \textbf{fiable}.
    \item[\textcolor{primary}{\faCheck}] Pour \textbf{comparer} les résultats entre laboratoires.
    \item[\textcolor{primary}{\faCheck}] Pour \textbf{prouver} la qualité d'une analyse (normes ISO 17025, accréditation).
    \item[\textcolor{primary}{\faCheck}] Pour estimer l'\textbf{incertitude} d'une mesure.
\end{itemize}

\section*{\faKey\quad Les 2 grands mots à retenir}

\begin{center}
\renewcommand{\arraystretch}{1.6}
\begin{tabularx}{\linewidth}{|>{\columncolor{soft}}p{3.5cm}|X|}
\hline
\textbf{\textcolor{primary}{\Large Justesse}} & La moyenne de mes résultats est-elle \textbf{proche de la vraie valeur} ? \\ \hline
\textbf{\textcolor{primary}{\Large Fidélité}} & Mes résultats sont-ils \textbf{groupés} (peu dispersés) ? \\ \hline
\end{tabularx}
\end{center}

\vspace{0.2cm}

\begin{retenir}[title={\faBullseye\quad Justesse + Fidélité = Exactitude}]
L'\textbf{exactitude}, c'est les deux à la fois : des résultats \textbf{centrés} ET \textbf{groupés}.
\end{retenir}

\vspace{0.3cm}

% Schéma cible
\begin{center}
\begin{tikzpicture}[scale=0.85]
  % Cible 1 : Juste et fidèle
  \draw[thick, primary] (0,0) circle (1.2);
  \draw[thick, primary] (0,0) circle (0.8);
  \draw[thick, primary] (0,0) circle (0.4);
  \fill[green] (0,0) circle (0.08);
  \foreach \x in {-0.15,0,0.15} \foreach \y in {-0.15,0,0.15}
    \fill[green] (\x,\y) circle (0.06);
  \node[below] at (0,-1.5) {\textbf{\textcolor{green}{Juste + Fidèle}}};

  % Cible 2 : Juste mais pas fidèle
  \draw[thick, primary] (4,0) circle (1.2);
  \draw[thick, primary] (4,0) circle (0.8);
  \draw[thick, primary] (4,0) circle (0.4);
  \fill[amber] (4,0) circle (0.08);
  \foreach \x/\y in {-0.9/0.7, 0.8/-0.5, 0.6/0.9, -0.7/-0.6, 0.3/0.5, -0.4/0.3}
    \fill[amber] (4+\x, \y) circle (0.06);
  \node[below] at (4,-1.5) {\textbf{\textcolor{amber}{Juste, pas fidèle}}};

  % Cible 3 : Ni juste ni fidèle
  \draw[thick, primary] (8,0) circle (1.2);
  \draw[thick, primary] (8,0) circle (0.8);
  \draw[thick, primary] (8,0) circle (0.4);
  \fill[red] (8,0) circle (0.08);
  \foreach \x/\y in {0.9/0.9, 1.0/0.5, 0.7/1.0, 1.1/0.3, 0.8/0.7}
    \fill[red] (8+\x, \y) circle (0.06);
  \node[below] at (8,-1.5) {\textbf{\textcolor{red}{Ni juste ni fidèle}}};
\end{tikzpicture}
\end{center}

\section*{\faFlask\quad Les 2 types de fidélité}

\begin{center}
\renewcommand{\arraystretch}{1.5}
\begin{tabularx}{\linewidth}{|p{3.3cm}|X|p{4cm}|}
\hline
\rowcolor{soft}
\textbf{Type} & \textbf{Quand ?} & \textbf{Entre...} \\ \hline
\textbf{Répétabilité} & Même jour, même labo, même opérateur, même machine. & ...des mesures faites \textbf{très proches}. \\ \hline
\textbf{Reproductibilité} & Labos différents, opérateurs différents, matériels différents. & ...des mesures \textbf{très variées}. \\ \hline
\end{tabularx}
\end{center}

\vspace{0.2cm}

\begin{labo}[title={\faFlask\quad Image simple}]
\textbf{Répétabilité} = je mesure 5 fois le même échantillon ce matin → les résultats doivent être très proches.\\
\textbf{Reproductibilité} = 5 labos différents mesurent le même échantillon → les résultats doivent rester proches malgré les conditions différentes.
\end{labo}

\section*{\faCalculator\quad Comment on fait ? (en 4 étapes)}

\begin{enumerate}[leftmargin=1.2cm, itemsep=4pt]
    \item \textbf{Plusieurs labos} analysent les mêmes échantillons.
    \item On \textbf{élimine} les valeurs aberrantes (tests de Cochran et Grubbs).
    \item On calcule deux nombres :
    \begin{itemize}
        \item \(s_r\) = écart-type de \textbf{répétabilité} (variabilité intra-labo).
        \item \(s_R\) = écart-type de \textbf{reproductibilité} (variabilité inter-labos).
    \end{itemize}
    \item On en déduit les \textbf{limites} :
    \[
    r = 2{,}8 \times s_r \qquad R = 2{,}8 \times s_R
    \]
\end{enumerate}

\begin{attention}[title={\faExclamationTriangle\quad À retenir absolument}]
\textbf{Deux mesures ne doivent pas différer de plus de \(r\) (même labo) ou \(R\) (labos différents)} dans 95\,\% des cas. Si c'est le cas → la mesure est douteuse.
\end{attention}

\section*{\faBuilding\quad Où on l'utilise en laboratoire ?}

\begin{center}
\renewcommand{\arraystretch}{1.4}
\begin{tabularx}{\linewidth}{|p{1.2cm}|X|}
\hline
\rowcolor{soft} \textbf{\faCheck} & \textbf{Utilisation} \\ \hline
\faCheckCircle & \textbf{Valider} une nouvelle méthode (HPLC, spectro...). \\ \hline
\faCheckCircle & \textbf{Comparer} les résultats entre laboratoires. \\ \hline
\faCheckCircle & \textbf{Estimer l'incertitude} d'une mesure. \\ \hline
\faCheckCircle & \textbf{Contrôler la qualité} au quotidien (stabilité). \\ \hline
\faCheckCircle & \textbf{Obtenir l'accréditation} ISO 17025. \\ \hline
\end{tabularx}
\end{center}

\section*{\faListOl\quad Les 6 parties de la norme (version express)}

\begin{center}
\renewcommand{\arraystretch}{1.4}
\begin{tabularx}{\linewidth}{|c|p{4cm}|X|}
\hline
\rowcolor{soft} \textbf{N°} & \textbf{Titre} & \textbf{En clair} \\ \hline
\textbf{1} & Principes et définitions & Le vocabulaire de base. \\ \hline
\textbf{2} & Méthode de base & Comment calculer répétabilité et reproductibilité. \\ \hline
\textbf{3} & Fidélité intermédiaire & Cas entre les deux (même labo, conditions variées). \\ \hline
\textbf{4} & Détermination de la justesse & Comment calculer le \textbf{biais}. \\ \hline
\textbf{5} & Méthodes alternatives & Méthodes \textbf{robustes} (sans éliminer les valeurs). \\ \hline
\textbf{6} & Utilisation pratique & Comment s'en servir au quotidien. \\ \hline
\end{tabularx}
\end{center}

\begin{retenir}[title={\faGraduationCap\quad Ce que tu dois savoir pour l'examen}]
\begin{enumerate}[leftmargin=1.2cm, itemsep=2pt]
    \item \textbf{Définir} justesse, fidélité, répétabilité, reproductibilité.
    \item \textbf{Distinguer} répétabilité (même labo) et reproductibilité (labos différents).
    \item \textbf{Écrire} les formules : \(r = 2{,}8\,s_r\) et \(R = 2{,}8\,s_R\).
    \item \textbf{Expliquer} à quoi sert l'ISO 5725 en laboratoire.
    \item \textbf{Citer} ISO/IEC 17025 comme norme qui renvoie à ISO 5725.
\end{enumerate}
\end{retenir}

\vspace{0.4cm}

\begin{center}
\textit{Fiche de synthèse — ISO 5725 — BTS BioALC — À garder dans ton classeur !}
\end{center}

\end{document}